\section{引言}
\label{sec:introduction}

从专家示范中进行模仿学习已成为获取复杂机器人操作技能的主流范式~\cite{chi2023diffusion,zhao2023act,brohan2023rt2}。在生成式方法中，\emph{流匹配}~\cite{lipman2023flow,tong2024improving}作为扩散策略的替代方案近年来受到广泛关注。通过学习一个确定性速度场，将简单先验分布通过常微分方程（ODE）传输到目标动作分布，流匹配具有多项优势：无需噪声调度的稳定训练、快速确定性采样，以及结合最优传输（OT-CFM）~\cite{tong2024improving}时的直线传输路径。VITA框架~\cite{gao2026vita}是这一方法的典型代表，证明了从视觉观测隐变量直接流向动作隐变量可以获得优于条件扩散策略的多模态覆盖能力。

尽管取得了上述进展，当前的流匹配策略——以及大多数模仿学习架构——在视觉编码阶段存在一个重要的局限性：

\paragraph{空间信息丢失。}标准观测编码器对卷积特征图施加全局平均池化，将所有空间结构压缩为单个向量。这丢弃了关于任务相关物体\emph{在视野中位置}的关键信息。对于焊缝跟踪等精密任务，策略不仅需要知道它\emph{看到了什么}，还需要空间布局——焊缝位置、接近方向以及工具与工件之间的相对几何关系。全局池化将一个$512 \times H \times W$的丰富特征图压缩为512维向量，空间分辨率从数百个位置降为零。

在本文中，我们提出\textbf{\fasafull{}（\fasa{}）}，一个通过帧间特征时间对比自监督发现任务相关空间区域的即插即用空间注意力模块（图~\ref{fig:architecture}）：

\begin{itemize}[leftmargin=*,itemsep=4pt]
    \item \textbf{自监督空间注意力}（第~\ref{sec:spatial}节）：\fasa{}的核心洞察是\textbf{连续帧之间变化最显著的区域对应于当前活跃的交互区域}。通过检测帧间特征变化，\fasa{}将空间特征分解为动态（交互）流和静态（上下文）流，编码变化质心和方向，并提取空间关键点坐标——提供比全局池化丰富$2\times$的空间信息，无需任何外部教师模型。

    \item \textbf{语义速度编码}（第~\ref{sec:temporal_diff}节）：与简单堆叠连续帧不同，我们计算显式的特征空间差分$\Delta \bm{f}_t = \bm{f}_t - \bm{f}_{t-1}$，编码特征空间中的方向性变化。这种"语义速度"在接近阶段大（场景快速变化）、跟踪阶段小（细微偏移），提供了隐式的阶段信号。
\end{itemize}

\fasa{}是\emph{完全自监督的}：它不需要空间标注或训练时的外部教师模型。\fasa{}的变化精炼网络通过策略的动作预测损失端到端训练，自动学习抑制无关变化（如光照）并放大任务相关变化（如工具-工件接触）。

我们将\fasa{}集成到VITA流匹配架构中，在一个具有挑战性的6自由度S型焊缝跟踪任务上进行评估，其中安川MR2010机械臂需要在仿真环境中接近随机放置的工件并精确跟随弯曲的焊接轨迹。实验表明，相比基础VITA策略和多个基线方法，\fasa{}显著提升了跟踪精度。

\paragraph{贡献。}本文的主要贡献如下：
\begin{enumerate}[leftmargin=*,itemsep=2pt]
    \item 我们识别出空间信息丢失是流匹配视觉运动策略的重要局限性，并提出\fasa{}（Flow-Aware Spatial Attention）作为即插即用的自监督解决方案。
    \item \fasa{}通过多统计量帧间特征时间对比自监督发现任务相关区域，将空间特征分解为动态流和静态流，编码变化质心和方向，并提取空间关键点坐标，消除了对外部教师模型的依赖。
    \item 结合特征差分的"语义速度"编码，\fasa{}仅需连续观测帧即可在标准卷积特征图上操作，使其广泛适用于本文研究的特定架构和任务之外。
    \item 我们在精密机器人跟踪任务上验证了\fasa{}的有效性，全面的消融研究证实了空间感知的关键贡献。
\end{enumerate}
